\section{Experiments}\label{sec:experiments}

The theorem concerns a ratio of expected risks, not the number of mistakes
in one graph. We therefore ask whether the empirical average
$\Mis(\widehat z,z)/R_{\mathrm{oracle}}$ remains on a constant scale as
$n$ increases. We compare only the complete spectral likelihood-ratio
clustering procedure of Section~\ref{restore:algorithm} with its
specified oracle classifier. There is
no comparison with other clustering methods, no fitted multiplicative
rescaling of the oracle curve, and no oracle initialization.

\subsection{Models, implementation, and the oracle denominator}

We use $n=1000,2000,4000,8000,16000,32000$, with 100 independent
Bernoulli graphs at every size in each of three settings. The two-community
model has equal community sizes and
\[
 P=p\begin{pmatrix}1&0.25\\0.25&1\end{pmatrix}.
\]
We first take $np=4\log n$. The three-community model has proportions
$(0.40,0.35,0.25)$ and
\[
 P=p\begin{pmatrix}
 1&0.10&0.15\\
 0.10&0.90&0.12\\
 0.15&0.12&0.80
 \end{pmatrix},\qquad np=6\log n.
\]
This second model has unequal expected degrees as well as unequal sizes.
As a third setting, we use the two-community matrix with
$np=6\sqrt{\log n}$. All matrix shapes and community proportions stay
fixed as $n$ changes. All three sequences satisfy $np\to\infty$ and
$np^2\to0$. The third sequence is $o(\log n)$ asymptotically;
on the displayed finite grid, however, $np$ still exceeds $\log n$.
It is not a numerical test of $np<\log n$ on these graph sizes.

The algorithm constants are fixed throughout: $b_0=3$, $c_0=0.01$,
and at most ten complete block-fit and simultaneous label-update rounds.
We stop earlier only if the labels are unchanged after a completed round.
The radius-count initializer uses the stated radii $1/8$ and $1/4$,
without $k$-means substitution or label-based restarts. The encoder retains
the $K$ eigenpairs of largest absolute eigenvalue, with relative eigensolver
tolerance $10^{-10}$. After encoding, the decoder receives only the original
degrees, the retained eigenpairs, and the fixed algorithm constants;
it has no access to adjacency entries, true labels, or true probabilities.
Block counts and symmetric block fits are computed from low-rank factors,
with the exact degree correction and all self-exclusions retained.

The denominator is not an observed error fraction from a single oracle
run, nor the asymptotic approximation $e^{-I}$. For a target in community
$k$, we use the independent count variables $b_{r}$ and the
contrasts $T_{k\ell}$ defined in Section~\ref{sec:statements}.
We sum their joint binomial probabilities over the error event in
\eqref{eq:oracle-risk}, including the full nonedge term, the $n_k-1$
own-community trials, and the fixed tie ordering, and then weight
community $k$ by $n_k/n$. This is a numerical evaluation of the specified
$R_{\mathrm{oracle}}$. Truncating the binomial supports discards total
probability less than $10^{-13}$ in each setting; no Poisson or Gaussian
approximation is used. The oracle errors observed on the simulated graphs
are also retained as an independent calibration check.

For each setting we divide the mean of the 100 misclassification
proportions by this fixed oracle risk. Pointwise 95\% Monte Carlo intervals
are Student-$t$ intervals for the mean across independent graphs, divided
by the same denominator. Vertices in one graph are not treated as
independent trials. These intervals describe simulation uncertainty,
not a simultaneous confidence bound over all $n$ or all model parameters.

\subsection{Risk ratios as the graph size grows}

Table~\ref{tab:experiment-risks} reports both risks and their ratio;
Figure~\ref{fig:experiment-ratios} shows the ratio without any vertical
rescaling. The balanced two-community ratios range from 2.84 to 2.96,
the unequal-size three-community ratios from 2.15 to 2.49, and the
more slowly growing two-community ratios from 3.19 to 3.31. In none
of these three sequences is there a sustained increase in the ratio
over the 32-fold increase in $n$.

In the first setting the mean misclassification proportion decreases
from 0.00791 to 0.00123 while the oracle risk decreases from 0.00271
to 0.000433. In the three-community setting the corresponding changes
are from 0.00404 to 0.000426 and from 0.00188 to 0.000198.
Thus the comparison concerns two decreasing risks, rather than a
constant error level at a fixed average degree. The ratios remain
above one; these results do not suggest convergence of the ratio to one.

\begin{table}[htbp]
\centering\small
\caption{Risk estimates for spectral likelihood-ratio clustering from 100 independent graphs per entry. The denominator is the raw-edge oracle risk computed from the exact Bernoulli block-count law. The last column gives a pointwise 95\% Monte Carlo confidence interval for the ratio, treating graphs, not vertices, as independent replicates.}\label{tab:experiment-risks}
\begin{tabular}{rcccc}
\toprule
$n$ & $R_{\mathrm{oracle}}$ & Mean $\Mis$ & Risk ratio & 95\% interval\\
\midrule
\multicolumn{5}{l}{\textit{Balanced $K=2$, $np=4\log n$}}\\[2pt]
1,000 & $2.709\times10^{-3}$ & $7.910\times10^{-3}$ & 2.92 & [2.69, 3.15]\\
2,000 & $1.893\times10^{-3}$ & $5.610\times10^{-3}$ & 2.96 & [2.79, 3.14]\\
4,000 & $1.315\times10^{-3}$ & $3.805\times10^{-3}$ & 2.89 & [2.75, 3.04]\\
8,000 & $9.094\times10^{-4}$ & $2.654\times10^{-3}$ & 2.92 & [2.78, 3.05]\\
16,000 & $6.280\times10^{-4}$ & $1.806\times10^{-3}$ & 2.88 & [2.78, 2.97]\\
32,000 & $4.334\times10^{-4}$ & $1.230\times10^{-3}$ & 2.84 & [2.74, 2.94]\\
\midrule
\multicolumn{5}{l}{\textit{Unequal-size $K=3$, $np=6\log n$}}\\[2pt]
1,000 & $1.881\times10^{-3}$ & $4.040\times10^{-3}$ & 2.15 & [1.92, 2.38]\\
2,000 & $1.228\times10^{-3}$ & $3.055\times10^{-3}$ & 2.49 & [2.31, 2.67]\\
4,000 & $7.981\times10^{-4}$ & $1.735\times10^{-3}$ & 2.17 & [2.01, 2.34]\\
8,000 & $4.913\times10^{-4}$ & $1.077\times10^{-3}$ & 2.19 & [2.04, 2.35]\\
16,000 & $3.119\times10^{-4}$ & $6.706\times10^{-4}$ & 2.15 & [2.02, 2.28]\\
32,000 & $1.977\times10^{-4}$ & $4.259\times10^{-4}$ & 2.15 & [2.04, 2.27]\\
\midrule
\multicolumn{5}{l}{\textit{Balanced $K=2$, $np=6\sqrt{\log n}$}}\\[2pt]
1,000 & $1.495\times10^{-2}$ & $4.956\times10^{-2}$ & 3.31 & [3.19, 3.44]\\
2,000 & $1.349\times10^{-2}$ & $4.422\times10^{-2}$ & 3.28 & [3.19, 3.36]\\
4,000 & $1.219\times10^{-2}$ & $4.018\times10^{-2}$ & 3.29 & [3.24, 3.35]\\
8,000 & $1.105\times10^{-2}$ & $3.609\times10^{-2}$ & 3.27 & [3.22, 3.31]\\
16,000 & $1.005\times10^{-2}$ & $3.246\times10^{-2}$ & 3.23 & [3.20, 3.26]\\
32,000 & $9.160\times10^{-3}$ & $2.922\times10^{-2}$ & 3.19 & [3.17, 3.21]\\
\bottomrule\end{tabular}\end{table}

\begin{figure}[htbp]
 \centering
 \includegraphics[width=.94\linewidth]{figures/risk_ratios.pdf}
 \caption{Mean misclassification divided by the exact raw-edge oracle risk.
 Each point uses 100 independent graphs. Error bars are the pointwise
 95\% Monte Carlo intervals in Table~\ref{tab:experiment-risks}.
 Stable finite-grid ratios support, but do not prove, the asymptotic
 constant-factor guarantee.}\label{fig:experiment-ratios}
\end{figure}

Across these 1800 primary graphs, every run reached an unchanged-label
update within ten rounds; none invoked the fixed fallback or activated
probability clipping. The largest individual misclassification proportions
were 0.015, 0.012, and 0.075 in the three settings, respectively. The
third setting included 78 isolated vertices in total, which were processed
without aborting classification of the other vertices. No degree cap was
active in these primary runs. Consequently, these performance measurements
exercise the full decoder, including degree correction, but do not by
themselves assess the statistical effect of active regularization.

\subsection{Implementation checks and finite-sample failures}

To distinguish numerical mistakes from poor finite-sample behavior,
we checked factorized counts, exact row sums, zero diagonals, symmetric
block fits, and likelihood scores against direct dense calculations
in 60 small-graph cases. These included 20 cases with active caps,
37 with retained negative eigenvalues, and 12 with a deliberately
isolated vertex. The largest absolute discrepancy was
$7.2\times10^{-15}$. Exhaustive enumeration of two small incident-row
experiments also agreed with the block-count oracle computation to
$8.4\times10^{-17}$. The largest relative eigen-equation residual
in the primary simulations was $5.1\times10^{-11}$.
These are implementation checks, not independent proofs of the risk theorem.

The choice of the stronger three-community sequence was made after a
separate pilot exposed unstable small-sample behavior at $np=4\log n$.
The pilot used different seeds from the primary experiment, and all its
outcomes are archived. We then ran a separate, independent stress test
with 100 graphs at each of $n=1000,2000,4000,8000$, retaining the same
three-community matrix and proportions but using $np=4\log n$.
Table~\ref{tab:stress} includes every run, not only successful ones.

\begin{table}[htbp]
\centering\small
\caption{Independent finite-sample stress test: the same three-community matrix and proportions, but $np=4\log n$. Each row includes all 100 simulated graphs. A large error means $\Mis>0.10$.}\label{tab:stress}
\begin{tabular}{rcccc}
\toprule
$n$ & Mean $\Mis$ & Risk ratio & 95\% interval & Large-error graphs\\
\midrule
1,000 & 0.10293 & 9.67 & [8.31, 11.02] & 43/100\\
2,000 & 0.06712 & 8.97 & [7.23, 10.71] & 9/100\\
4,000 & 0.02346 & 4.37 & [3.63, 5.11] & 1/100\\
8,000 & 0.01133 & 2.89 & [2.81, 2.96] & 0/100\\
\bottomrule\end{tabular}\end{table}

The weakest three-community samples can therefore fail substantially:
at $n=1000$, 43 of 100 graphs have more than 10\% misclassification,
and the largest error is 0.41. In 59 of those 100 graphs the retained
three-dimensional space includes a negative eigenvalue, even though
all three nonzero eigenvalues of the completed population matrix are
positive. One recorded failure retains eigenvalues approximately
$(12.86,9.37,-7.30)$ and selects two initial representatives from the
same true community; the true labels are used only for this post-run
diagnosis. This exposes a finite-sample noise-direction and initialization
failure, rather than an omitted nonedge term or an incorrect degree
correction. It is not asserted to explain every failed run. At $n=8000$
the same weaker sequence has mean risk ratio 2.89 and no graph with
more than 10\% error.

No graph was discarded, rerun with an oracle-selected seed, or repaired
using its true labels. The code, frozen designs, all 2200 non-pilot
replicate records, all 24 pilot records, oracle calculations, and
implementation checks accompany the source. These experiments provide
finite-sample evidence for constant-scale oracle comparisons on the
three specified sequences, while documenting a range in which the
method is not yet reliable. They do not establish a uniform numerical
constant for the entire parameter class or replace the asymptotic theorem.
